\subsection{1. Mô tả chi tiết bài toán}
\subsubsection{1.1. Bài toán phát hiện chữ bị làm giả trong tài liệu}
\indent\indent Nghiên cứu tập trung giải quyết bài toán phát hiện và phân loại các vùng văn bản bị làm giả (Tampered Text Detection) trên các ảnh tài liệu số hóa. Khác với bài toán nhận dạng chữ quang học (OCR) thông thường chỉ đọc nội dung ký tự, bài toán pháp y tài liệu đòi hỏi mô hình phải phân tích sâu các đặc trưng vi mô ở cấp độ pixel để xác định tính xác thực của từng ký tự hoặc khối văn bản.\vspace{0.3cm}

Các hình thức làm giả văn bản phổ biến trong tài liệu bao gồm:\vspace{-0.1cm}
\begin{itemize}[label={--}, itemsep=1pt]
    \item \textbf{Sao chép - Dán (Copy-Move):} Kẻ giả mạo sao chép một cụm ký tự hoặc con số từ một vị trí trong cùng tài liệu (ví dụ: con số 0 trong ngày tháng) và dán đè vào vị trí khác (như số tiền giao dịch). Dấu hiệu đặc trưng là sự trùng khớp tuyệt đối về độ nhiễu và hoa văn font chữ nhưng lại có viền tiếp giáp bất thường.
    \item \textbf{Ghép nối văn bản (Text Splicing):} Lấy chữ số hoặc từ ngữ từ một tài liệu nguồn khác ghép vào tài liệu đích. Trường hợp này thường để lại các dấu vết sai lệch rõ rệt về độ phân giải, độ sắc nét viền, góc nghiêng hoặc bảng màu nền.
    \item \textbf{Chỉnh sửa hoặc thay thế phông chữ (Digital Font Replacement):} Sử dụng công cụ đồ họa máy tính để gõ lại nội dung chữ mới đè lên nền đã xóa. Dấu hiệu là sự thiếu tự nhiên về độ răng cưa (aliasing), sai khác về khoảng cách ký tự (kerning) và độ tương phản quang sai.
    \item \textbf{Xóa văn bản (Text Erasure / Inpainting):} Sử dụng thuật toán lấp đầy nội dung để xóa bỏ chữ ký, con dấu hoặc các điều khoản quan trọng.
\end{itemize}

\subsubsection{1.2. Mục tiêu hệ thống và yêu cầu suy luận}
\indent\indent Để một hệ thống giám định tài liệu có giá trị sử dụng cao trong thực tế, hệ thống không chỉ cần độ chính xác cao mà còn phải thỏa mãn hai tiêu chí thiết yếu:\vspace{-0.1cm}
\begin{itemize}[label={--}, itemsep=1pt]
    \item \textbf{Tính gọn nhẹ và thời gian thực:} Mô hình phân loại cần có kích thước nhỏ, tốc độ suy luận nhanh để có thể triển khai trên các máy chủ xác thực tài liệu hàng loạt hoặc tích hợp vào thiết bị di động của giao dịch viên.
    \item \textbf{Khả năng giải thích suy luận (Explainability):} Thay vì chỉ trả về một nhãn phân loại hay điểm xác suất trừu tượng, hệ thống cần cung cấp được lý do cụ thể vì sao vùng chữ đó bị coi là giả mạo (ví dụ: ``\textit{Phát hiện viền cắt không tự nhiên xung quanh cụm số 500,000, font chữ có độ sắc nét không tương đồng với các ký tự lân cận}'').
\end{itemize}

Nhằm đáp ứng trọn vẹn các yêu cầu này, đề tài đề xuất mô hình phối hợp giữa kiến trúc phân loại tích chập tối ưu **EfficientNet-B0** và mô hình thị giác - ngôn ngữ **SmolVLM2-500M** chuyên biệt cho tác vụ trích xuất suy luận logic.
